\chapter{Chapter09 — Gaussian Quadrant Asymptotics}
%%% generated by the blueprint pipeline from FABL.Chapter09.GaussianQuadrantAsymptotics
\section{Overview}

Declarations of the Lean module \texttt{FABL.Chapter09.GaussianQuadrantAsymptotics}.

\section{Declarations}

\begin{definition}[Positive correlation parameter]
\lean{FABL.positiveCorrelation}
\label{FABL.positiveCorrelation}
\leanok
For $\rho\in(0,1)$, define its positive correlation parameter in $(-1,1)$ to be the same
real number $\rho$.
\end{definition}

\begin{definition}[Correlation residual variance]
\lean{FABL.correlationResidualVariance}
\statementsource{mossel-neeman-gaussian-noise-stability}{
  pdf-b53c1a5aa27d-p012-b005,
  pdf-b53c1a5aa27d-p012-b006
}
\label{FABL.correlationResidualVariance}
\leanok
For a real correlation parameter $\rho$ with $0<\rho<1$, define the residual variance to
be $1-\rho^2$.
\end{definition}

\begin{theorem}[Residual variance formula]
\lean{FABL.coe_correlationResidualVariance}
\label{FABL.coe_correlationResidualVariance}
\leanok
\uses{FABL.correlationResidualVariance}
\difficulty{1}
For every real number $\rho$ with $0<\rho<1$, the real value of the nonnegative residual
variance $1-\rho^2$ is exactly $1-\rho^2$.
\end{theorem}

\begin{theorem}[Positive residual variance]
\lean{FABL.correlationResidualVariance_pos}
\label{FABL.correlationResidualVariance_pos}
\leanok
\uses{FABL.correlationResidualVariance}
\difficulty{2}
For every real correlation parameter $\rho$ with $0<\rho<1$, the residual variance
$1-\rho^2$ is positive.
\end{theorem}

\begin{theorem}[Nonzero residual variance]
\lean{FABL.correlationResidualVariance_ne_zero}
\label{FABL.correlationResidualVariance_ne_zero}
\leanok
\uses{FABL.correlationResidualVariance, FABL.correlationResidualVariance_pos}
\difficulty{1}
For every real correlation parameter $\rho$ with $0<\rho<1$, the residual variance
$1-\rho^2$ is nonzero.
\end{theorem}

\begin{definition}[Correlated Gaussian density]
\lean{FABL.correlatedGaussianDensity}
\statementsource{mossel-neeman-gaussian-noise-stability}{
  pdf-b53c1a5aa27d-p003-b002,
  pdf-b53c1a5aa27d-p012-b005
}
\label{FABL.correlatedGaussianDensity}
\leanok
\uses{FABL.correlationResidualVariance}
For a real correlation parameter $\rho$ with $0<\rho<1$ and real numbers $z_1,z_2$,
define the correlated Gaussian density at $(z_1,z_2)$ as the standard Gaussian density
at $z_1$ multiplied by the Gaussian density at $z_2$ with mean $\rho z_1$ and variance
$1-\rho^2$.
\end{definition}

\begin{theorem}[Nonnegativity of the correlated Gaussian density]
\lean{FABL.correlatedGaussianDensity_nonneg}
\label{FABL.correlatedGaussianDensity_nonneg}
\leanok
\uses{FABL.correlatedGaussianDensity}
\difficulty{1}
For every $0<\rho<1$ and $z_1,z_2\in\mathbb{R}$, the correlated Gaussian density at
$(z_1,z_2)$ is nonnegative.
\end{theorem}

\begin{theorem}[Positivity of the correlated Gaussian density]
\lean{FABL.correlatedGaussianDensity_pos}
\label{FABL.correlatedGaussianDensity_pos}
\leanok
\uses{FABL.correlatedGaussianDensity, FABL.correlationResidualVariance, FABL.correlationResidualVariance_ne_zero}
\difficulty{1}
For every $0<\rho<1$ and $z_1,z_2\in\mathbb{R}$, the correlated Gaussian density at
$(z_1,z_2)$ is strictly positive.
\end{theorem}

\begin{theorem}[Derivative of a Gaussian density]
\lean{FABL.hasDerivAt_gaussianPDFReal_posVariance}
\label{FABL.hasDerivAt_gaussianPDFReal_posVariance}
\leanok
\difficulty{5}
Let $\mu,x\in\mathbb R$ and let $v>0$ be a variance. The Gaussian density with
mean $\mu$ and variance $v$ is differentiable at $ x$, with derivative
$-\frac{x-\mu}{v}$ times its value at $x$.
\end{theorem}

\begin{theorem}[Gaussian density vanishes at positive infinity]
\lean{FABL.tendsto_gaussianPDFReal_posVariance_atTop}
\label{FABL.tendsto_gaussianPDFReal_posVariance_atTop}
\leanok
\difficulty{4}
For every real mean $\mu$ and every positive variance $v$, the density at $x$ of the
Gaussian distribution with mean $\mu$ and variance $v$ tends to $0$ as $x$ tends to
positive infinity.
\end{theorem}

\begin{theorem}[Integrability of a centered Gaussian density]
\lean{FABL.integrable_centered_mul_gaussianPDFReal}
\label{FABL.integrable_centered_mul_gaussianPDFReal}
\leanok
\difficulty{4}
Let $\mu\in\mathbb{R}$ and let $v\ge 0$ be a nonzero real parameter. The function
$x\mapsto (x-\mu)g_{\mu,v}(x)$ is integrable on $\mathbb{R}$, where $g_{\mu,v}$ is the
Gaussian density with mean $\mu$ and parameter $v$.
\end{theorem}

\begin{theorem}[Gaussian tail integral identity]
\lean{FABL.integral_Ioi_centered_mul_gaussianPDFReal}
\label{FABL.integral_Ioi_centered_mul_gaussianPDFReal}
\leanok
\uses{FABL.hasDerivAt_gaussianPDFReal_posVariance, FABL.integrable_centered_mul_gaussianPDFReal, FABL.tendsto_gaussianPDFReal_posVariance_atTop}
\difficulty{4}
Let $\mu,t\in\mathbb{R}$ and let $v>0$. If $g_{\mu,v}$ denotes the Gaussian density with
mean $\mu$ and parameter $v$, then
\[\int_t^\infty (x-\mu)g_{\mu,v}(x)\,dx=v\,g_{\mu,v}(t).\]
\end{theorem}

\begin{theorem}[Formula for the correlated Gaussian density]
\lean{FABL.correlatedGaussianDensity_formula}
\label{FABL.correlatedGaussianDensity_formula}
\leanok
\uses{FABL.coe_correlationResidualVariance, FABL.correlatedGaussianDensity, FABL.correlationResidualVariance_pos}
\difficulty{3}
Let $0<\rho<1$ and $z_1,z_2\in\mathbb{R}$. The correlated Gaussian density is
$\displaystyle \frac{1}{\sqrt{2\pi}\sqrt{2\pi(1-\rho^2)}}\exp\!\left(-\frac{z_1^2-2\rho
z_1z_2+z_2^2}{2(1-\rho^2)}\right)$.
\end{theorem}

\begin{theorem}[Symmetry of the correlated Gaussian density]
\lean{FABL.correlatedGaussianDensity_comm}
\label{FABL.correlatedGaussianDensity_comm}
\leanok
\uses{FABL.correlatedGaussianDensity, FABL.correlatedGaussianDensity_formula}
\difficulty{3}
Let $0<\rho<1$ and $z_1,z_2\in\mathbb{R}$. The joint density obtained by taking a
standard normal first coordinate and, conditional on its value $z_1$, a normal second
coordinate with mean $\rho z_1$ and variance $1-\rho^2$ is symmetric in its arguments:
its value at $(z_1,z_2)$ equals its value at $(z_2,z_1)$.
\end{theorem}

\begin{theorem}[A centered correlated Gaussian tail integral]
\lean{FABL.integral_Ioi_second_centered_mul_correlatedGaussianDensity}
\proofsource{Analysis-of-Boolean-Functions-by-Ryan-ODonnell}{pdf-ffbe9c7b3306-p274-b004}
\label{FABL.integral_Ioi_second_centered_mul_correlatedGaussianDensity}
\leanok
\uses{FABL.correlatedGaussianDensity, FABL.correlationResidualVariance, FABL.correlationResidualVariance_pos, FABL.integral_Ioi_centered_mul_gaussianPDFReal}
\difficulty{4}
Let $0<\rho<1$ and let $z_1,t\in\mathbb R$. For the density $g_\rho(z_1,z_2)$ of a pair
of standard Gaussian variables with correlation $\rho$, defined as the standard Gaussian
density at $z_1$ times the Gaussian density at $z_2$ with mean $\rho z_1$ and variance
$1-\rho^2$, one has
\[\int_{z_2>t}(z_2-\rho z_1)g_\rho(z_1,z_2)\,dz_2=(1-\rho^2)g_\rho(z_1,t).\]
\end{theorem}

\begin{theorem}[Integral ioi first centered mul correlated gaussian density]
\lean{FABL.integral_Ioi_first_centered_mul_correlatedGaussianDensity}
\statementsource{Analysis-of-Boolean-Functions-by-Ryan-ODonnell}{pdf-ffbe9c7b3306-p274-b004}
\label{FABL.integral_Ioi_first_centered_mul_correlatedGaussianDensity}
\leanok
\uses{FABL.correlatedGaussianDensity, FABL.correlatedGaussianDensity_comm, FABL.correlationResidualVariance, FABL.correlationResidualVariance_pos, FABL.integral_Ioi_centered_mul_gaussianPDFReal}
\difficulty{3}
% dataset-pipeline: stub -- prose pending build_informal_statements.py
Statement pending regeneration.
\end{theorem}

\begin{theorem}[Weighted quadrant integral and density]
\lean{FABL.gaussian_weighted_quadrant_integral_eq_variance_sq_mul_density}
\proofsource{Analysis-of-Boolean-Functions-by-Ryan-ODonnell}{pdf-ffbe9c7b3306-p274-b004}
\label{FABL.gaussian_weighted_quadrant_integral_eq_variance_sq_mul_density}
\leanok
\uses{FABL.correlatedGaussianDensity, FABL.integral_Ioi_first_centered_mul_correlatedGaussianDensity, FABL.integral_Ioi_second_centered_mul_correlatedGaussianDensity}
\difficulty{4}
Let $0<\rho<1$ and $t\in\mathbb R$. The iterated integral over $z_1>t$ and $z_2>t$ of
$(z_2-\rho z_1)(z_1-\rho t)$ times the density of a centered bivariate standard normal
pair with correlation $\rho$ equals $(1-\rho^2)^2$ times that density evaluated at
$(t,t)$.
\end{theorem}

\begin{theorem}[Diagonal correlated Gaussian density identity]
\lean{FABL.variance_sq_mul_correlatedGaussianDensity_diagonal}
\label{FABL.variance_sq_mul_correlatedGaussianDensity_diagonal}
\leanok
\uses{FABL.correlatedGaussianDensity, FABL.correlatedGaussianDensity_formula, FABL.correlationResidualVariance_pos}
\difficulty{5}
For $0<\rho<1$ and $t\in\bbr$, the correlated Gaussian density on the diagonal satisfies
\[(1-\rho^2)^2\,q_\rho(t,t)=\frac{(\sqrt{1-\rho^2})^3}{2\pi}\exp\!\left(-\frac{t^2}{1+\rho}\right),\]
where $q_\rho$ denotes the density of a pair of standard normal variables with
correlation $\rho$.
\end{theorem}

\begin{theorem}[Weighted Gaussian quadrant integral]
\lean{FABL.gaussian_weighted_quadrant_integral}
\proofsource{Analysis-of-Boolean-Functions-by-Ryan-ODonnell}{pdf-ffbe9c7b3306-p274-b004}
\label{FABL.gaussian_weighted_quadrant_integral}
\leanok
\uses{FABL.correlatedGaussianDensity, FABL.gaussian_weighted_quadrant_integral_eq_variance_sq_mul_density, FABL.variance_sq_mul_correlatedGaussianDensity_diagonal}
\difficulty{1}
Let $0<\rho<1$ and $t\in\mathbb R$. For the density of a centered bivariate standard
normal pair with correlation $\rho$, the iterated integral over $z_1>t$ and $z_2>t$ of
$(z_2-\rho z_1)(z_1-\rho t)$ times that density equals
$\frac{(1-\rho^2)^{3/2}}{2\pi}\exp\!\left(-\frac{t^2}{1+\rho}\right)$.
\end{theorem}

\begin{definition}[Gaussian quadrant probability at a threshold]
\lean{FABL.gaussianQuadrantAtThreshold}
\statementsource{Analysis-of-Boolean-Functions-by-Ryan-ODonnell}{
  pdf-ffbe9c7b3306-p138-b001,
  pdf-ffbe9c7b3306-p274-b002
}
\statementsource{mossel-neeman-gaussian-noise-stability}{
  pdf-b53c1a5aa27d-p003-b004,
  pdf-b53c1a5aa27d-p012-b003
}
\label{FABL.gaussianQuadrantAtThreshold}
\leanok
\uses{FABL.gaussianQuadrantProbability, FABL.positiveCorrelation, FABL.standardGaussianUpperTailOpen}
For a real correlation parameter $\rho$ with $0<\rho<1$ and a real threshold $t$, define
the Gaussian quadrant probability at threshold $t$ to be the probability, for a
positively correlated standard Gaussian pair with correlation parameter $\rho$, that
both coordinates lie in the open upper tail above $t$.
\end{definition}

\begin{theorem}[Gaussian quadrant probability at a threshold]
\lean{FABL.gaussianQuadrantAtThreshold_eq_measure}
\label{FABL.gaussianQuadrantAtThreshold_eq_measure}
\leanok
\uses{FABL.CorrelationPlane, FABL.correlatedGaussianMeasure, FABL.gaussianQuadrantAtThreshold, FABL.gaussianQuadrantProbability, FABL.gaussianUpperRightQuadrant, FABL.standardGaussianUpperQuantile_upperTail}
\difficulty{3}
Let $0<\rho<1$ and $t\in\mathbb R$. For a pair of standard Gaussian random variables
with correlation $\rho$, the probability that both coordinates exceed $t$ equals the
correlated Gaussian measure of the upper right quadrant above $t$.
\end{theorem}

\begin{theorem}[Gaussian quadrant at threshold le tail rpow]
\lean{FABL.gaussianQuadrantAtThreshold_le_tail_rpow}
\statementsource{Analysis-of-Boolean-Functions-by-Ryan-ODonnell}{
  pdf-ffbe9c7b3306-p274-b004,
  pdf-ffbe9c7b3306-p274-b006
}
\label{FABL.gaussianQuadrantAtThreshold_le_tail_rpow}
\leanok
\uses{FABL.gaussianQuadrantAtThreshold, FABL.hammingUpperTailIndicator, FABL.mean, FABL.noiseStability, FABL.positiveCorrelation, FABL.smallSetExpansion, FABL.standardGaussianUpperTail, FABL.standardGaussianUpperTail_pos, FABL.tendsto_expect_hammingUpperTailIndicator, FABL.tendsto_noiseStability_hammingUpperTailIndicator}
\difficulty{5}
% dataset-pipeline: stub -- prose pending build_informal_statements.py
Statement pending regeneration.
\end{theorem}

\begin{definition}[Gaussian conditional slope]
\lean{FABL.gaussianConditionalSlope}
\statementsource{Analysis-of-Boolean-Functions-by-Ryan-ODonnell}{pdf-ffbe9c7b3306-p274-b003}
\statementsource{mossel-neeman-gaussian-noise-stability}{
  pdf-b53c1a5aa27d-p012-b004,
  pdf-b53c1a5aa27d-p012-b005
}
\label{FABL.gaussianConditionalSlope}
\leanok
For a real parameter $0<\rho<1$, define the Gaussian conditional slope by $\displaystyle
\frac{1-\rho}{\sqrt{1-\rho^2}}$.
\end{definition}

\begin{theorem}[Positivity of the Gaussian conditional slope]
\lean{FABL.gaussianConditionalSlope_pos}
\label{FABL.gaussianConditionalSlope_pos}
\leanok
\uses{FABL.correlationResidualVariance_pos, FABL.gaussianConditionalSlope}
\difficulty{3}
For every $
ho\in(0,1)$, the Gaussian conditional slope $rac{1-
ho}{\sqrt{1-
ho^2}}$ is
positive.
\end{theorem}

\begin{theorem}[Gaussian conditional slope identity]
\lean{FABL.sqrt_residual_mul_gaussianConditionalSlope}
\label{FABL.sqrt_residual_mul_gaussianConditionalSlope}
\leanok
\uses{FABL.correlationResidualVariance_pos, FABL.gaussianConditionalSlope}
\difficulty{3}
For every real number $\rho$ with $0<\rho<1$, define the Gaussian conditional slope by
$c(\rho)=(1-\rho)/\sqrt{1-\rho^2}$. Then $\sqrt{1-\rho^2}\,c(\rho)=1-\rho$.
\end{theorem}

\begin{theorem}[Gaussian quadrant at threshold ge rectangle]
\lean{FABL.gaussianQuadrantAtThreshold_ge_rectangle}
\statementsource{Analysis-of-Boolean-Functions-by-Ryan-ODonnell}{
  pdf-ffbe9c7b3306-p274-b004,
  pdf-ffbe9c7b3306-p274-b006
}
\label{FABL.gaussianQuadrantAtThreshold_ge_rectangle}
\leanok
\uses{FABL.continuous_gaussianCorrelationMap, FABL.correlationResidualVariance_pos, FABL.gaussianConditionalSlope, FABL.gaussianCorrelationMap, FABL.gaussianQuadrantAtThreshold, FABL.gaussianQuadrantAtThreshold_eq_measure, FABL.gaussianUpperRightQuadrant, FABL.isOpen_gaussianUpperRightQuadrant, FABL.sqrt_residual_mul_gaussianConditionalSlope, FABL.standardGaussianUpperTail, FABL.standardGaussianUpperTail_eq_measureReal_Ioi}
\difficulty{5}
% dataset-pipeline: stub -- prose pending build_informal_statements.py
Statement pending regeneration.
\end{theorem}

\begin{theorem}[Gaussian measure on an open interval]
\lean{FABL.standardGaussian_measureReal_Ioo_eq_integral}
\label{FABL.standardGaussian_measureReal_Ioo_eq_integral}
\leanok
\difficulty{3}
For all real numbers $a$ and $b$, the standard Gaussian measure of the open interval
$(a,b)$ equals the integral of its density over that interval:
$\gamma((a,b))=\int_{(a,b)}\varphi(x)\,dx$, where $\varphi$ is the standard Gaussian
density.
\end{theorem}

\begin{theorem}[Monotonicity of the Gaussian density on the nonnegative axis]
\lean{FABL.standardGaussianDensity_antitone_of_nonneg}
\label{FABL.standardGaussianDensity_antitone_of_nonneg}
\leanok
\difficulty{3}
For real numbers $x$ and $y$ with $0\le x\le y$, the standard Gaussian density at $y$ is
at most its value at $x$.
\end{theorem}

\begin{theorem}[Gaussian density bound on an interval]
\lean{FABL.density_endpoint_div_le_standardGaussian_interval}
\label{FABL.density_endpoint_div_le_standardGaussian_interval}
\leanok
\uses{FABL.standardGaussianDensity_antitone_of_nonneg, FABL.standardGaussian_measureReal_Ioo_eq_integral}
\difficulty{4}
For every real number $t>0$, the standard Gaussian density at $t+1/t$, divided by $t$,
is at most the standard Gaussian measure of the open interval $(t,t+1/t)$.
\end{theorem}

\begin{theorem}[Shifted Gaussian density bound]
\lean{FABL.exp_neg_three_halves_mul_density_le_shifted_density}
\label{FABL.exp_neg_three_halves_mul_density_le_shifted_density}
\leanok
\difficulty{5}
For every real $t\ge 1$, the standard Gaussian density satisfies
$e^{-3/2}\,\varphi(t)\le \varphi(t+1/t)$, where $\varphi$ denotes the density of a
Gaussian with mean $0$ and variance $1$.
\end{theorem}

\begin{theorem}[A Gaussian interval probability lower bound]
\lean{FABL.exp_neg_three_halves_mul_density_div_le_standardGaussian_interval}
\label{FABL.exp_neg_three_halves_mul_density_div_le_standardGaussian_interval}
\leanok
\uses{FABL.density_endpoint_div_le_standardGaussian_interval, FABL.exp_neg_three_halves_mul_density_le_shifted_density}
\difficulty{3}
For every real $t\ge1$, the standard Gaussian density at $t$, multiplied by
$e^{-3/2}/t$, is at most the standard Gaussian probability of the interval $(t,t+1/t)$.
\end{theorem}

\begin{theorem}[Square of the Gaussian conditional slope]
\lean{FABL.gaussianConditionalSlope_sq}
\label{FABL.gaussianConditionalSlope_sq}
\leanok
\uses{FABL.correlationResidualVariance_pos, FABL.gaussianConditionalSlope}
\difficulty{4}
For every $
ho\in(0,1)$, the square of the Gaussian conditional slope satisfies
$igl(rac{1-
ho}{\sqrt{1-
ho^2}}igr)^2=rac{1-
ho}{1+
ho}$.
\end{theorem}

\begin{theorem}[Squared Gaussian conditional slope identity]
\lean{FABL.one_add_gaussianConditionalSlope_sq}
\label{FABL.one_add_gaussianConditionalSlope_sq}
\leanok
\uses{FABL.gaussianConditionalSlope, FABL.gaussianConditionalSlope_sq}
\difficulty{3}
For every real number $\rho$ with $0<\rho<1$, if $s=(1-\rho)/\sqrt{1-\rho^2}$ is the
Gaussian conditional slope, then $1+s^2=\frac{2}{1+\rho}$.
\end{theorem}

\begin{theorem}[Gaussian quadrant at threshold lower mills]
\lean{FABL.gaussianQuadrantAtThreshold_lower_mills}
\statementsource{Analysis-of-Boolean-Functions-by-Ryan-ODonnell}{
  pdf-ffbe9c7b3306-p274-b003,
  pdf-ffbe9c7b3306-p274-b004,
  pdf-ffbe9c7b3306-p274-b006
}
\label{FABL.gaussianQuadrantAtThreshold_lower_mills}
\leanok
\uses{FABL.density_mul_t_div_one_add_sq_le_standardGaussianUpperTail, FABL.exp_neg_three_halves_mul_density_div_le_standardGaussian_interval, FABL.gaussianConditionalSlope, FABL.gaussianConditionalSlope_pos, FABL.gaussianQuadrantAtThreshold, FABL.gaussianQuadrantAtThreshold_ge_rectangle, FABL.standardGaussianUpperTail}
\difficulty{4}
% dataset-pipeline: stub -- prose pending build_informal_statements.py
Statement pending regeneration.
\end{theorem}

\begin{theorem}[Product of Gaussian densities along a conditional slope]
\lean{FABL.standardGaussianDensity_mul_slope_density}
\label{FABL.standardGaussianDensity_mul_slope_density}
\leanok
\uses{FABL.gaussianConditionalSlope, FABL.gaussianConditionalSlope_sq}
\difficulty{4}
Let $0<\rho<1$ and $t\in\mathbb{R}$, and set $s=(1-\rho)/\sqrt{1-\rho^2}$. The product
of the standard Gaussian densities at $t$ and $st$ equals
$(2\pi)^{-1}\exp\!\left(-t^2/(1+\rho)\right)$.
\end{theorem}

\begin{theorem}[Conditional slope denominator comparison]
\lean{FABL.gaussianConditionalSlope_div_large_denominator_le}
\label{FABL.gaussianConditionalSlope_div_large_denominator_le}
\leanok
\uses{FABL.gaussianConditionalSlope, FABL.gaussianConditionalSlope_pos}
\difficulty{4}
Let $
ho\in(0,1)$ and $t\inbR$ with $1\le t$. If $s=(1-
ho)/\sqrt{1-
ho^2}$ is the
Gaussian conditional slope, then $rac{s}{(1+s^2)t^2}\lerac{s}{1+s^2t^2}$.
\end{theorem}

\begin{definition}[Gaussian quadrant lower constant]
\lean{FABL.gaussianQuadrantLowerConstant}
\statementsource{Analysis-of-Boolean-Functions-by-Ryan-ODonnell}{pdf-ffbe9c7b3306-p274-b002}
\label{FABL.gaussianQuadrantLowerConstant}
\leanok
\uses{FABL.gaussianConditionalSlope}
For a real parameter $0<\rho<1$, define the Gaussian quadrant lower constant to be
$\displaystyle \frac{e^{-3/2}}{2\pi}\,\frac{s}{1+s^2}$, where $\displaystyle
s=\frac{1-\rho}{\sqrt{1-\rho^2}}$.
\end{definition}

\begin{theorem}[Positivity of the Gaussian quadrant lower constant]
\lean{FABL.gaussianQuadrantLowerConstant_pos}
\label{FABL.gaussianQuadrantLowerConstant_pos}
\leanok
\uses{FABL.gaussianConditionalSlope_pos, FABL.gaussianQuadrantLowerConstant}
\difficulty{3}
For every correlation parameter $0<\rho<1$, the Gaussian quadrant lower constant is
positive.
\end{theorem}

\begin{theorem}[Gaussian quadrant at threshold lower exp div sq]
\lean{FABL.gaussianQuadrantAtThreshold_lower_exp_div_sq}
\statementsource{Analysis-of-Boolean-Functions-by-Ryan-ODonnell}{
  pdf-ffbe9c7b3306-p274-b003,
  pdf-ffbe9c7b3306-p274-b004,
  pdf-ffbe9c7b3306-p274-b006
}
\label{FABL.gaussianQuadrantAtThreshold_lower_exp_div_sq}
\leanok
\uses{FABL.gaussianConditionalSlope, FABL.gaussianConditionalSlope_div_large_denominator_le, FABL.gaussianQuadrantAtThreshold, FABL.gaussianQuadrantAtThreshold_lower_mills, FABL.gaussianQuadrantLowerConstant, FABL.standardGaussianDensity_mul_slope_density}
\difficulty{5}
% dataset-pipeline: stub -- prose pending build_informal_statements.py
Statement pending regeneration.
\end{theorem}

\begin{definition}[Soft power exponent at zero]
\lean{FABL.HasSoftPowerAtZero}
\label{FABL.HasSoftPowerAtZero}
\leanok
Let $F\colon(0,1)\to\mathbb R$ and $a\in\mathbb R$. Say that $F$ has soft power exponent
$a$ at zero when $\log(F(\mu))/\log(\mu)$ tends to $a$ as $\mu$ tends to zero through
positive values.
\end{definition}

\begin{theorem}[Tendsto log standard gaussian upper tail div sq]
\lean{FABL.tendsto_log_standardGaussianUpperTail_div_sq}
\statementsource{Analysis-of-Boolean-Functions-by-Ryan-ODonnell}{pdf-ffbe9c7b3306-p274-b003}
\label{FABL.tendsto_log_standardGaussianUpperTail_div_sq}
\leanok
\uses{FABL.log_inv_standardGaussianUpperTail_isEquivalent_sq, FABL.standardGaussianUpperTail}
\difficulty{5}
% dataset-pipeline: stub -- prose pending build_informal_statements.py
Statement pending regeneration.
\end{theorem}

\begin{theorem}[Logarithm divided by a square]
\lean{FABL.tendsto_log_div_sq_atTop}
\label{FABL.tendsto_log_div_sq_atTop}
\leanok
\difficulty{4}
As the real variable $t$ tends to $+\infty$, $rac{\\log t}{t^2}$ tends to $0$.
\end{theorem}

\begin{theorem}[Decay of a reciprocal quadratic]
\lean{FABL.tendsto_const_div_sq_atTop}
\label{FABL.tendsto_const_div_sq_atTop}
\leanok
\difficulty{3}
For every real constant $c$, the function $t\mapsto c/t^2$ tends to $0$ as
$t\to+\infty$.
\end{theorem}

\begin{theorem}[Logarithmic Gaussian scale identity]
\lean{FABL.log_gaussianQuadrantLowerScale}
\label{FABL.log_gaussianQuadrantLowerScale}
\leanok
\uses{FABL.gaussianQuadrantLowerConstant, FABL.gaussianQuadrantLowerConstant_pos}
\difficulty{3}
For $0<\rho<1$, define $s=(1-\rho)/\sqrt{1-\rho^2}$ and
$C=e^{-3/2}s/\bigl((1+s^2)2\pi\bigr)$. For every $t>0$,
\[\log\!\left(C\frac{e^{-t^2/(1+\rho)}}{t^2}\right)=\log C-\frac{t^2}{1+\rho}-2\log t.\]
\end{theorem}

\begin{theorem}[Tendsto log gaussian quadrant lower scale div tail log]
\lean{FABL.tendsto_log_gaussianQuadrantLowerScale_div_tailLog}
\statementsource{Analysis-of-Boolean-Functions-by-Ryan-ODonnell}{
  pdf-ffbe9c7b3306-p274-b002,
  pdf-ffbe9c7b3306-p274-b003,
  pdf-ffbe9c7b3306-p274-b004
}
\label{FABL.tendsto_log_gaussianQuadrantLowerScale_div_tailLog}
\leanok
\uses{FABL.gaussianQuadrantLowerConstant, FABL.log_gaussianQuadrantLowerScale, FABL.standardGaussianUpperTail, FABL.standardGaussianUpperTail_lt_one, FABL.standardGaussianUpperTail_pos, FABL.tendsto_const_div_sq_atTop, FABL.tendsto_log_div_sq_atTop, FABL.tendsto_log_standardGaussianUpperTail_div_sq}
\difficulty{5}
% dataset-pipeline: stub -- prose pending build_informal_statements.py
Statement pending regeneration.
\end{theorem}

\begin{theorem}[Tendsto log gaussian quadrant at threshold div tail log]
\lean{FABL.tendsto_log_gaussianQuadrantAtThreshold_div_tailLog}
\statementsource{Analysis-of-Boolean-Functions-by-Ryan-ODonnell}{
  pdf-ffbe9c7b3306-p274-b002,
  pdf-ffbe9c7b3306-p274-b003,
  pdf-ffbe9c7b3306-p274-b004,
  pdf-ffbe9c7b3306-p274-b006
}
\label{FABL.tendsto_log_gaussianQuadrantAtThreshold_div_tailLog}
\leanok
\uses{FABL.gaussianQuadrantAtThreshold, FABL.gaussianQuadrantAtThreshold_le_tail_rpow, FABL.gaussianQuadrantAtThreshold_lower_exp_div_sq, FABL.gaussianQuadrantLowerConstant, FABL.gaussianQuadrantLowerConstant_pos, FABL.standardGaussianUpperTail, FABL.standardGaussianUpperTail_lt_one, FABL.standardGaussianUpperTail_pos, FABL.tendsto_log_gaussianQuadrantLowerScale_div_tailLog}
\difficulty{5}
% dataset-pipeline: stub -- prose pending build_informal_statements.py
Statement pending regeneration.
\end{theorem}

\begin{theorem}[Soft power of the Gaussian quadrant probability]
\lean{FABL.gaussianQuadrantProbability_hasSoftPower}
\proofsource{Analysis-of-Boolean-Functions-by-Ryan-ODonnell}{pdf-ffbe9c7b3306-p274-b002}
\label{FABL.gaussianQuadrantProbability_hasSoftPower}
\leanok
\uses{FABL.HasSoftPowerAtZero, FABL.gaussianQuadrantAtThreshold, FABL.gaussianQuadrantProbability, FABL.positiveCorrelation, FABL.standardGaussianUpperQuantile, FABL.standardGaussianUpperTailOpen, FABL.standardGaussianUpperTail_quantile, FABL.tendsto_log_gaussianQuadrantAtThreshold_div_tailLog, FABL.tendsto_standardGaussianUpperQuantile_atBot}
\difficulty{4}
For every correlation parameter $0<\rho<1$, the Gaussian quadrant probability as a
function of its upper-tail parameter has soft power $2/(1+\rho)$ at zero: as
$\mu\to0^+$, the ratio of the logarithm of that probability at $\mu$ to $\log\mu$ tends
to $2/(1+\rho)$.
\end{theorem}

\begin{theorem}[A normalized Gaussian density identity]
\lean{FABL.normalized_mills_power_eq_exp_div_sq}
\label{FABL.normalized_mills_power_eq_exp_div_sq}
\leanok
\difficulty{5}
For every $\rho\in(0,1)$ and every real $t>0$, the following identity holds, where
$\varphi(t)$ is the density of the standard normal distribution: $\displaystyle
\frac{(\sqrt{2\pi})^{2/(1+\rho)}\bigl(\varphi(t)/t\bigr)^{2/(1+\rho)}}{t^{2-2/(1+\rho)}}=\frac{\exp\bigl(-t^2/(1+\rho)\bigr)}{t^2}$.
\end{theorem}

\begin{definition}[Gaussian quadrant polylogarithmic constant]
\lean{FABL.gaussianQuadrantPolylogConstant}
\label{FABL.gaussianQuadrantPolylogConstant}
\leanok
\uses{FABL.gaussianQuadrantLowerConstant}
For $\rho\in(0,1)$, define the Gaussian quadrant polylogarithmic constant as
$c(\rho)\,\sqrt{2\pi}^{\,2/(1+\rho)}$, where $c(\rho)$ is the Gaussian quadrant lower
constant at correlation $\rho$.
\end{definition}

\begin{theorem}[Positivity of the Gaussian quadrant polylogarithmic constant]
\lean{FABL.gaussianQuadrantPolylogConstant_pos}
\label{FABL.gaussianQuadrantPolylogConstant_pos}
\leanok
\uses{FABL.gaussianQuadrantLowerConstant_pos, FABL.gaussianQuadrantPolylogConstant}
\difficulty{1}
For every correlation parameter $0<\rho<1$, the Gaussian quadrant polylogarithmic
constant is positive.
\end{theorem}

\begin{theorem}[Lower bound for Gaussian quadrant probability]
\lean{FABL.gaussianQuadrantAtThreshold_lower_tail_rpow_div_quantilePower}
\label{FABL.gaussianQuadrantAtThreshold_lower_tail_rpow_div_quantilePower}
\leanok
\uses{FABL.gaussianQuadrantAtThreshold, FABL.gaussianQuadrantAtThreshold_lower_exp_div_sq, FABL.gaussianQuadrantLowerConstant, FABL.gaussianQuadrantPolylogConstant, FABL.gaussianQuadrantPolylogConstant_pos, FABL.normalized_mills_power_eq_exp_div_sq, FABL.standardGaussianUpperTail, FABL.standardGaussianUpperTail_le_density_div, FABL.standardGaussianUpperTail_pos}
\difficulty{4}
Let $0<\rho<1$ and $t\geq 1$. The Gaussian quadrant probability at threshold $t$ is at
least the product of the Gaussian quadrant polylogarithmic constant, the standard
Gaussian upper tail at $t$ raised to $2/(1+\rho)$, and the reciprocal of $t$ raised to
$2-2/(1+\rho)$.
\end{theorem}

\begin{theorem}[Gaussian quadrant probability bounds]
\lean{FABL.eventually_gaussianQuadrantProbability_polylog_bounds}
\label{FABL.eventually_gaussianQuadrantProbability_polylog_bounds}
\leanok
\uses{FABL.gaussianQuadrantAtThreshold, FABL.gaussianQuadrantAtThreshold_le_tail_rpow, FABL.gaussianQuadrantAtThreshold_lower_tail_rpow_div_quantilePower, FABL.gaussianQuadrantPolylogConstant, FABL.gaussianQuadrantProbability, FABL.positiveCorrelation, FABL.standardGaussianUpperQuantile, FABL.standardGaussianUpperTailOpen, FABL.standardGaussianUpperTail_quantile, FABL.tendsto_standardGaussianUpperQuantile_atBot}
\difficulty{4}
Fix a correlation parameter $\rho\in(0,1)$. For all sufficiently small $\mu\in(0,1)$,
the probability that two standard Gaussian coordinates with correlation $\rho$ both
exceed the standard Gaussian upper-tail quantile at mass $\mu$ lies between
$C_\rho\,\mu^{2/(1+\rho)}z_\mu^{2-2/(1+\rho)}$ and $\mu^{2/(1+\rho)}$, where $z_\mu$ is
that quantile and $C_\rho$ is the stated positive constant depending on $\rho$.
\end{theorem}

\begin{definition}[Gaussian tail at a shifted threshold]
\lean{FABL.gaussianBadTailAtThreshold}
\label{FABL.gaussianBadTailAtThreshold}
\leanok
\uses{FABL.standardGaussianUpperTail}
Let $Z$ be a standard normal random variable. For $
ho	ext{ with }0<
ho<1$ and $t	ext{
in }br$, define the Gaussian bad tail at threshold $t$ to be
$bP\bigl(Z>(t-1)/\rho\bigr)$.
\end{definition}

\begin{theorem}[Shifted correlation quotient tends to infinity]
\lean{FABL.tendsto_shift_div_correlation_atTop}
\label{FABL.tendsto_shift_div_correlation_atTop}
\leanok
\difficulty{3}
For every real $\rho$ with $0<\rho<1$, the quantity $(t-1)/\rho$ tends to $+\infty$ as
the real variable $t$ tends to $+\infty$.
\end{theorem}

\begin{theorem}[Squared shifted quotient limit]
\lean{FABL.tendsto_shift_div_correlation_sq_div_sq}
\label{FABL.tendsto_shift_div_correlation_sq_div_sq}
\leanok
\difficulty{4}
For every real $\rho$ with $0<\rho<1$, as $t$ tends to $+\infty$, $((t-1)/\rho)^2/t^2$
tends to $1/\rho^2$.
\end{theorem}

\begin{theorem}[Gaussian bad-tail logarithmic rate]
\lean{FABL.tendsto_log_gaussianBadTail_div_sq}
\label{FABL.tendsto_log_gaussianBadTail_div_sq}
\leanok
\uses{FABL.gaussianBadTailAtThreshold, FABL.standardGaussianUpperTail, FABL.tendsto_log_standardGaussianUpperTail_div_sq, FABL.tendsto_shift_div_correlation_atTop, FABL.tendsto_shift_div_correlation_sq_div_sq}
\difficulty{5}
Let $0<\rho<1$, let $Q(t)=\mathbb{P}(Z>t)$ for a standard normal random variable $Z$,
and set $B_\rho(t)=Q((t-1)/\rho)$. As $t\to+\infty$, $\frac{\log B_\rho(t)}{t^2}$ tends
to $-\frac{1}{2\rho^2}$.
\end{theorem}

\begin{theorem}[Ratio of Gaussian tail logarithms]
\lean{FABL.tendsto_log_gaussianBadTail_div_tailLog}
\label{FABL.tendsto_log_gaussianBadTail_div_tailLog}
\leanok
\uses{FABL.gaussianBadTailAtThreshold, FABL.standardGaussianUpperTail, FABL.standardGaussianUpperTail_lt_one, FABL.standardGaussianUpperTail_pos, FABL.tendsto_log_gaussianBadTail_div_sq, FABL.tendsto_log_standardGaussianUpperTail_div_sq}
\difficulty{4}
Let $0<\rho<1$, let $Q(t)=\mathbb{P}(Z>t)$ for a standard normal random variable $Z$,
and set $B_\rho(t)=Q((t-1)/\rho)$. As $t\to+\infty$, $\frac{\log B_\rho(t)}{\log Q(t)}$
tends to $\frac{1}{\rho^2}$.
\end{theorem}

\begin{definition}[Gaussian bad tail probability]
\lean{FABL.gaussianBadTailProbability}
\label{FABL.gaussianBadTailProbability}
\leanok
\uses{FABL.gaussianBadTailAtThreshold, FABL.standardGaussianUpperQuantile}
For $\rho,\mu\in(0,1)$, define the Gaussian bad tail probability to be the Gaussian bad
tail at threshold equal to the standard Gaussian upper quantile of $\mu$.
\end{definition}

\begin{theorem}[Soft power of the Gaussian bad-tail probability]
\lean{FABL.gaussianBadTailProbability_hasSoftPower}
\label{FABL.gaussianBadTailProbability_hasSoftPower}
\leanok
\uses{FABL.HasSoftPowerAtZero, FABL.gaussianBadTailProbability, FABL.standardGaussianUpperTail_quantile, FABL.tendsto_log_gaussianBadTail_div_tailLog, FABL.tendsto_standardGaussianUpperQuantile_atBot}
\difficulty{3}
For every $
ho\in(0,1)$, as $\mu	o0$ through $(0,1)$, the logarithmic ratio of the
Gaussian bad-tail probability to $\log\mu$ tends to $1/
ho^2$; that is, $rac{\log
P_
ho(\mu)}{\log\mu}	o1/
ho^2$, where $P_
ho(\mu)$ is the Gaussian bad-tail probability
defined using the standard Gaussian upper quantile and threshold $(t-1)/
ho$.
\end{theorem}

\begin{theorem}[Comparison of two powers]
\lean{FABL.mainPower_lt_badTailPower}
\label{FABL.mainPower_lt_badTailPower}
\leanok
\difficulty{4}
For every real number $\rho$ with $0<\rho<1$, one has
\[\frac{2}{1+\rho}<\frac{1}{\rho^2}.\]
\end{theorem}

\begin{theorem}[Little-o bound for the Gaussian bad-tail probability]
\lean{FABL.gaussianBadTailProbability_isLittleO_mainPower}
\label{FABL.gaussianBadTailProbability_isLittleO_mainPower}
\leanok
\uses{FABL.gaussianBadTailProbability, FABL.gaussianBadTailProbability_hasSoftPower, FABL.mainPower_lt_badTailPower, FABL.standardGaussianUpperTail_pos}
\difficulty{6}
For every $
ho\in(0,1)$, as $\mu	o0$ through $(0,1)$, the Gaussian bad-tail probability
$P_
ho(\mu)$ is little-o of $\mu^{2/(1+
ho)}$; equivalently,
$rac{P_
ho(\mu)}{\mu^{2/(1+
ho)}}	o0$.
\end{theorem}
